\begin{itemize}
\item \textbf{Polymerace:} monomery se spojují do řetězce, např. otevřením násobné vazby. Probíhá iniciace, růst a ukončení řetězce; nemění se chemické složení základní jednotky a nevzniká vedlejší nízkomolekulární produkt. Příklad: ethylen $\to$ polyethylen.
\item \textbf{Polyadice:} reagují funkční skupiny monomerů, vzniká chemicky odlišná základní jednotka, ale bez odštěpení malé molekuly. Příklad: diisokyanát a diol $\to$ polyuretan.
\item \textbf{Polykondenzace:} reakcí funkčních skupin vzniká polymer a současně se odštěpují malé molekuly, např. voda či HCl. Příklad: diamin a dikarboxylová kyselina $\to$ polyamid a voda.
\end{itemize}
